\section{Обобщённый нисходящий алгоритм для поиска путей с КС ограничениям}
\label{sec:CFPQ_GLL}
\tikzsetfigurename{CFPQ_GLL_}

GLL довольно естественно обобщается на граф~\cite{Grigorev:2017:CPQ:3166094.3166104}: позициями входа теперь будем считать не индексы линейного слова, а вершины графа. В самом же алгоритме требуется внести лишь два небольших дополнения:

\begin{enumerate}
  \item Теперь при обработке терминала <<следующих>> символов может быть несколько~--- рассматриваем каждый из них отдельно, сдвигаясь по соответствующему ребру. В результате, при одном чтении можем получить несколько новых дескрипторов, но они независимы, потому просто ставим их в рабочее множество $ R $.
  \item При обработке нетерминала, аналогично, правила управляющей таблицы применяются для каждого из <<следующих>> символов в графе. Соответственно новых дескрипторов будет сгенерировано больше, но все они по-прежнему независимы и просто добавляются в рабочее множество $ R $.
\end{enumerate}

Подробное описание алгоритма и псевдокод представлены в работе~\cite{Grigorev:2017:CPQ:3166094.3166104}. Существует ещё одно обобщение нисходящего синтаксического анализа для решения задачи КС достижимости~\cite{MEDEIROS201975}, которое предполагает непосредственное обобщение LL(k) алгоритма, что приводит к аналогичному результату, однако теряется связь с некоторыми построениями.

Основанный на нисходящем анализе алгоритма поиска путей с контекстно-свободными ограничениями имеет следующие особенности.
\begin{enumerate}
  \item Необходимо явно задавать начальную вершину, поэтому хорошо подходит для задач поиска путей с одним источником (single-source) или с небольшим количеством источников. Для поиска путей между всеми парами вершин необходимо явным образом все указать стартовыми.
  \item Является направленным сверху вниз~--- обходит граф последовательно начиная с указанной стартовой вершины и строит вывод, начиная со стартового нетерминала. Как следствие, в отличие от алгоритмов на основе линейной алгебры и Хеллингса, обойдёт только подграф, необходимый для построения ответа. В среднем это меньше, чем весь граф, который обрабатывается другими алгоритмами.
  \item Естественным образом строит множество путей в виде сжатого леса разбора.
  \item Использует существенно более тяжеловесные стуктуры данных и плохо распараллеливается (на практике). Как следствие, при решении задачи достижимости для всех пар путей проигрывает алгоритмам на основе линейной алгебры.
\end{enumerate}

Частным случаем применения задачи КС достижимости является синтаксический анализ с неоднозначной токенизацией, то есть ситуацией, когда несколько пересекающихся подстрок во входной строке символов могут задавать разные лексические единицы и не возможно сделать однозначный выбор на этапе лексического анализа.
Например, для строки \verb|x = a~---b| возможны несколько вариантов токенизации.
\begin{enumerate}
  \item \verb|ID(x) OP_EQ ID(a) OP_MINUS OP_DECRIMENT ID(b)|
  %\item \verb|ID(x) OP_EQ ID(A) OP_MINUS OP_UN_MINUS ID(b)|
  \item \verb|ID(x) OP_EQ ID(a) OP_DECRIMENT OP_MINUS ID(b)|
\end{enumerate}

В таком случае на вход синтаксическому анализатору можно подать DAG, содержащий все возможные варианты токенизации. Для нашего примера он может выглядеть следующим образом:

\begin{center}
        \begin{tikzpicture}[node distance=4cm,shorten >=1pt,on grid,auto]
    \node[state] (q_0) at (0,0)  {$0$};
    \node[state] (q_1) at (2.5,0)  {$1$};
    \node[state] (q_2) at (5.5,0)  {$2$};
    \node[state] (q_3) at (8,0)  {$3$};
    \node[state] (q_4) at (10,-2)  {$4$};
    \node[state] (q_5) at (11.5,0) {$6$};
    \node[state] (q_6) at (10,2) {$5$};
    \node[state] (q_7) at (14,0) {$7$};
    \path[->]
    (q_0) edge  node {ID(x)} (q_1)
    (q_1) edge  node {OP\_EQ} (q_2)
    (q_2) edge  node {ID(a)} (q_3)
    (q_3) edge[left]  node {OP\_MINUS} (q_4)
    (q_4) edge[right]  node {OP\_DECRIMENT} (q_5)
    (q_3) edge  node {OP\_DECRIMENT} (q_6)
    (q_6) edge node {OP\_MINUS} (q_5)
    (q_5) edge  node {ID(b)} (q_7);
    \end{tikzpicture}

\end{center}

Далее будем проверять наличие пути из старовой (нулевой) вершины в конечную (соответствующую концу строки). Если таких путей оказалось несколько, то нужны дополнительные средства для выбора нужного дерева разбора. Данная идея рассматривается в работе~\cite{10.1145/3357766.3359532}.

Напоследок сделаем небольшое замечание об эффективной реализации: в качестве рабочего множества $ R $ можно использовать несколько различных структур данных и, как правило, выбирают очередь. Однако иногда (в особенности для графов) лучше использовать стек дескрипторов, так как в этом случае выше локальность данных~--- мы кладём пачку дескрипторов, соответствующих исходящим рёбрам. И если граф представлен списком смежности, то исходящие будут храниться рядом и их лучше обработать сразу.

%\section{Вопросы и задачи}
%\begin{enumerate}
%  \item Проведите алгоритм GLL для грамматики $ S \to a S b S \mid \varepsilon$. Правда ли, что эта грамматика принадлежит классу $ LL(1) $? Пронаблюдайте, как использование GSS вырождается в работу с обычным стеком.
%  \item Доразберите все не рассмотренные в примере \ref{gll:example1} дескрипторы, постройте полный GSS.
%\end{enumerate}
